\TieudeT{PHONG TỎA VẬT LÍ 11 - DAO ĐỘNG VÀ SÓNG}
\subsubsection*{PHẦN I. Thí sinh trả lời câu hỏi từ câu 1 đến câu 18. Mỗi câu hỏi thí sinh chỉ chọn một phương án}
\setcounter{ex}{0}
	\begin{ex}
		Phát biểu nào sau đây đúng về dao động điều hoà?
		\choice
		{Quỹ đạo là đường hình sin}
		{\True Quỹ đạo là một đoạn thẳng}
		{Vận tốc tỉ lệ thuận với thời gian}
		{Gia tốc tỉ lệ thuận với thời gian}
		%\haidong{5}
	\end{ex}
	%\loai{Xác định li độ}
	\begin{ex}
		Một chất điểm dao động điều hoà với phương trình $x=5\cos(10\pi t+\dfrac{\pi}{3})\ cm$. Li độ của vật khi pha dao động bằng $\pi$ là	
		\choice
		{$5\ cm$}
		{\True $-5\ cm$}
		{$2,5\ cm$}
		{$-2,5\ cm$}
		\loigiai{
			Li độ của vật khi pha dao động bằng $\pi$ là $-5\ cm$.
		}
		%\haidong{5}
	\end{ex}
	\begin{ex}%[3P1Y2-1][Loai1]
		Trong quá trình dao động, gia tốc của vật có độ lớn cực tiểu khi vật
		\choice
		{\True đi qua vị trí cân bằng}
		{ở biên (dương hoặc âm)}
		{ở biên âm ($x = - A$)}
		{ở biên dương ($x = A$)}
		\loigiai{
			
		}
	\end{ex}
	%\loai{Cho hai gia tốc. Tìm gia tốc tại vị trí thứ 3}
	\begin{ex} %[3P1-2.2K][Loai1]  
		Một chất điểm đang dao động điều hòa trên một đoạn thẳng. Gia tốc của chất điểm khi ngang qua vị trí M và N lần lượt là $a_M = 30\ cm/s^2$ và $a_N = 70\ cm/s^2$. Khi đi ngang qua trung điểm I của đoạn MN, gia tốc chuyển động của chất điểm là
		\choice
		{$30\ cm/s^2$}
		{$40\ cm/s^2$}
		{\True $50\ cm/s^2$}
		{$26\ cm/s^2$}
		\loigiai{
			\begin{itemize}
				\item Ta có: $\begin{cases}
					a_M = -\omega^2.x_M = 30\ cm/s^2.\\ 
					a_N = -\omega^2.x_N = 70\ cm/s^2.\end{cases}$
				\item I là trung điểm của đoạn MN thì $x_I = \dfrac{x_M + x_N}{2}$\\
				$\Rightarrow a_I = -\omega^2.\dfrac{x_M + x_N}{2} = \dfrac{a_M + a_N}{2} = \dfrac{30 + 70}{2} = 50\ cm/s^2$.
		\end{itemize}}
		%\haidong{15}
	\end{ex}
	%\loai{Cho pha, trạng thái dao động tìm đại lượng khác}
	\begin{ex}%[3P1B4-1][Loai1]
		Một chất điểm dao động điều hoà theo hàm cosin với chu kì $2\ s$ và có vận tốc $- 1\ m/s$ vào lúc pha dao động bằng $\dfrac{\pi}{4}\ rad$. Biên độ dao động là
		\choice
		{$15\ cm$}
		{\True $45\ cm$}
		{$0,25\ m$}
		{$35\ cm$}
		\loigiai{
			\begin{itemize}
				\item Ta có: $\omega  = \pi$  rad/s.
				\item Khi $\phi =\dfrac{\pi}{4}:\ x =\dfrac{A\sqrt{2}}{2}\ (-) \Rightarrow v=-\dfrac{v_{1\max}\sqrt{2}}{2}=-\dfrac{\omega A\sqrt{2}}{2}=-1\Rightarrow A=45\ cm$.
			\end{itemize}
		}
		%\haidong{15}
	\end{ex}
	\begin{ex}%[3P1-8.2GK][Loai1] 
		Cho một chất điểm dao động điều hòa theo phương trình vận tốc $v = \pi\cos\left(\dfrac{\pi}{2}t + \pi\right)$, với v tính bằng $cm/s$ và t tính bằng giây. Tính từ lúc $t=0$, vật đi qua vị trí li độ $x = -\sqrt{2}\ cm$ lần thứ $6$ tại thời điểm
		\choice
		{$5,5\ s$}
		{$19\ s$}
		{\True $9,5\ s$}
		{$1,5\ s$}
		\loigiai{
			\immini{\begin{itemize}
					\item Ta có: $v = \pi.cos\left(\dfrac{\pi}{2}.t+ \pi\right)\ cm$/s $ \Rightarrow x = 2cos\left(\dfrac{\pi}{2}.t+ \dfrac{\pi}{2}\right)\ cm$.
					\item $T = 4$ s. Ta biểu diễn vị trí các điểm có li độ $-\sqrt{2}\ cm$ trên đường tròn lượng giác như hình vẽ.
					\item Ta có: $6 = 2.2 + 2\Rightarrow t = 2.T+ \Delta t$.
					\item Từ hình vẽ ta thấy: $\Delta \varphi = \dfrac{\pi}{2} + \dfrac{\pi}{4} = \dfrac{3\pi}{4} \Rightarrow \Delta t = \dfrac{3T}{8}\Rightarrow t = 2T+ \dfrac{3T}{8} = \dfrac{19T}{8} = 9,5\ s$.
			\end{itemize}}{\begin{tikzpicture}[scale=1,thick,>=stealth']
					\tkzInit[ymin=-3,ymax=3,xmin=-3,xmax=3]\tkzClip
					\tkzDefPoints{-2.5/0/m, 2.5/0/n, 0/0/o, 0/2.5/p, 0/-2.5/q}
					\tkzDefShiftPoint[o](225:2){2}
					\tkzDefShiftPoint[o](135:2){1}
					\tkzDefShiftPoint[o](90:2){0}
					\tkzDefPointBy[projection=onto m--n](1) \tkzGetPoint{h}
					\draw(o)circle(2cm);
					\tkzDrawSegments[dashed](1,2)
					\tkzDrawSegments(p,q)
					\tkzDrawSegments[->](m,n)
					\tkzDrawSegments[blue,->](o,1 o,2 o,0)
					\tkzDrawPoints[size=3](o,1,h,0,2)
					\tkzMarkAngles[size=0.7cm,arrows=->](0,o,2)
					\node at (0)[above right]{$M_0$};
					\node at (1)[above left]{$M_1$};
					\node at (2)[below left]{$M_2$};
					\node at (2,0)[below right]{$2$};
					\node at (h)[below]{$-\sqrt{2}$};
			\end{tikzpicture}}
		}
		%\haidong{15}
		
	\end{ex}
	\begin{ex} %[3P1B3-1][Loai1]  
		Trong con lắc lò xo nếu ta tăng khối lượng vật nặng lên $4$ lần và độ cứng tăng $2$ lần thì tần số dao động của vật
		\choice
		{tăng $2$ lần}
		{giảm $2$ lần}
		{tăng $\sqrt{2}$ lần}
		{\True giảm $\sqrt{2}$ lần}
		\loigiai{
			\begin{itemize}
				\item Ta có: $f_1 = \dfrac{1}{2\pi} \sqrt{\dfrac{k}{m}}$. Khi tăng m và k: $f_2 = \dfrac{1}{2\pi} \sqrt{\dfrac{2k}{4m}} = \dfrac{1}{2\pi} \sqrt{\dfrac{k}{m}}.\dfrac{1}{\sqrt{2}} = \dfrac{f_1}{\sqrt{2}}$.
		\end{itemize}}
		%\haidong{15}
	\end{ex}
	%\loai{Số dao động trong thời gian $\Delta t$}
	\begin{ex}%[3P1B3-2]
		Một con lắc lò xo dao động điều hoà. Trong khoảng thời gian  $\Delta t$, con lắc thực hiện $60$ dao động toàn phần. Thay đổi khối lượng con lắc một lượng $440\ g$ thì cũng trong khoảng thời gian $\Delta t$ ấy, nó thực hiện $50$ dao động toàn phần. Khối lượng ban đầu của con lắc là
		\choice
		{$1,44\ kg$}
		{$0,6\ kg$}
		{$0,8\ kg$}
		{\True $1\ kg$}
		\loigiai{
			Ta có: $\begin{cases}
				\dfrac{\Delta t}{60} = 2\pi \sqrt{\dfrac{m}{k}} \\
				\dfrac{\Delta t}{50}=2\pi \sqrt{\dfrac{m\pm 0,44}{k}} 
			\end{cases} 
			\Rightarrow \dfrac{5}{6} = \sqrt{\dfrac{m}{m \pm 0,44}} \Rightarrow \dfrac{5}{6} = \sqrt{\dfrac{m}{m + 0,44}} \Rightarrow m = 1\ kg.$
		}
		%\haidong{15}
	\end{ex}
	
	\begin{ex}%[3P2-8.2K][Loai1]
		Dây căng ngang hai đầu cố định với chiều dài $\ell $, trên dây có sóng dừng. Nếu tăng chiều dài của dây lên gấp đôi (hai đầu vẫn cố định) thì dây có $10$ bụng sóng, nếu tăng chiều dài thêm $30\ cm$ (hai đầu vẫn cố định) thì trên dây có $8$ nút sóng. Biết tần số, tốc độ sóng trên dây không đổi trong quá trình thay đổi chiều dài dây. Chiều dài ban đầu $\ell$ của dây là
		\choice
		{$50\ cm$}
		{\True $75\ cm$}
		{$150\ cm$}
		{$100\ cm$}
		\loigiai{
			\begin{itemize}
				\item Ban đầu: $f=n\dfrac{v}{2\ell }$.
				.			\item Tăng chiều dài gấp đôi: $f=10\dfrac{v}{4\ell }$.
				\item Tăng chiều dài thêm 30 cm: $f=7\dfrac{v}{2\left( \ell +30 \right)} \Rightarrow \ell  = 75$ cm.
				\item 
				\item 
			\end{itemize}
		}
	\end{ex}
	
	
	%\loai{Viết phương trình dao động}
	\begin{ex}%[3P2-3.1B][Loai1]
		Một sóng hình sin truyền theo chiều dương của trục Ox với phương trình dao động của nguồn sóng (đặt tại O) là $u_O = 4\cos100\pi t\ cm$. Ở điểm M (theo hướng Ox) cách O một phần tư bước sóng, phần tử môi trường dao động với phương trình là
		\choice
		{$u_M = 4\cos(100\pi t + \pi)\ cm$}
		{$u_M = 4\cos100\pi t\ cm$}
		{\True $u_M = 4\cos(100\pi t - 0,5\pi)\ cm$}
		{$u_M = 4\cos(100\pi t + 0,5\pi)\ cm$}
		\loigiai{
			\begin{itemize}
				\item Sóng từ O đến M nên M chậm pha hơn so với O góc $\dfrac{2\pi d}{\lambda }=\dfrac{\pi }{2} \Rightarrow  u_M = 4\cos(100\pi t - 0,5\pi )$ cm.
				\item 
				\item 
			\end{itemize}
		}
	\end{ex}
	\loai{Viết phương trình dao động tại 1 điểm - ngược chiều truyền}
	\begin{ex}%[3P2-3.1K][Loai1]
		Sóng cơ truyền từ A đến B trên sợi dây AB rất dài với tốc độ $20\ m/s$. Tại điểm N trên dây cách A $75\ cm$, các phần tử ở đó dao động với phương trình $u_N = 3\cos20\pi t\ cm$, $t$ tính bằng $s$. Bỏ qua sự giảm biên độ. Phương trình dao động của phần tử tại điểm M trên dây cách A $50\ cm$ là
		\choice
		{\True $u_M = 3\cos(20\pi t + \dfrac{\pi}{4})\ cm$}
		{$u_M = 3\cos(20\pi t - \dfrac{\pi}{4})\ cm$}
		{$u_M = 3\cos(20\pi t + \dfrac{\pi}{2})\ cm$}
		{$u_M = 3\cos(20\pi t - \dfrac{\pi}{2})\ cm$}
		\loigiai{
			\begin{itemize}
				\item Rõ ràng M gần nguồn hơn nên nhanh pha hơn N góc $\dfrac{2\pi d}{\lambda }=\dfrac{2\pi \left( 75-25 \right)}{200}=\dfrac{\pi }{4} \Rightarrow  u_M = 3\cos(20\pi t + \dfrac{\pi }{4})\ cm$.
				\item 
				\item 
			\end{itemize}
		}
	\end{ex}
	\begin{ex}%[3P2-11.2K][Loai1]
		Sóng âm khi truyền trong chất rắn có thể là sóng dọc hoặc sóng ngang và lan truyền với tốc độ khác nhau. Tại trung tâm phòng chống thiên tai nhận được hai tín hiệu sóng từ một vụ động đất cách nhau một khoảng thời gian $240\ s$. Biết tốc độ truyền sóng ngang và tốc độ truyền sóng dọc trong lòng đất có giá trị lần lượt là $5\ km/s$ và $8\ km/s$. Tâm chấn động cách nơi nhận tín hiệu một khoảng 
		\choice
		{\True $3200\ km$}
		{$570\ km$}
		{$730\ km$}
		{$3500\ km$}
		\loigiai{
			\begin{itemize}
				\item Gọi d là khoảng cách từ tâm chấn động đến nơi nhận tín hiệu, ta có:
				$\dfrac{d}{5}-\dfrac{d}{8}=240\Rightarrow d=3200$ km.
				\item 
				\item 
			\end{itemize}
		}
	\end{ex}
	\begin{ex} %[3P4-4.2B][Loai4]
		Một sóng điện từ lan truyền trong chân không có bước sóng $3000\ m$. Lấy $c=3.10^8\ m/s$. Biết trong sóng điện từ, thành phần từ trường tại một điểm biến thiên điều hòa với chu kì T. Giá trị của T là
		\choice
		{$4.10^{-6}\ s$}
		{$2.10^{-5}\ s$}
		{\True $10^{-5}\ s$}
		{$3.10^{-6}\ s$}
		\loigiai{
			\begin{itemize}
				\item Chu kì của sóng: $T=\dfrac{\lambda}{c}=10^{-5}\ s$.
			\end{itemize}
		}
	\end{ex}
	\begin{ex} %[3P5-3.2K][Loai7]
		Trong thí nghiệm Young về giao thoa ánh sáng đơn sắc, khoảng cách giữa $4$ vân sáng liên tiếp trên màn quan sát là $2,4\ mm$. Khoảng vân trên màn là
		\choice
		{$1,6\ mm$}	
		{$1,2\ mm$}
		{$0,6\ mm$}	
		{\True $0,8\ mm$}
		\loigiai{
			\begin{itemize}
				\item 4 vân sáng: $3i = 2,4 \Rightarrow i = 0,8\ mm$.
				\item 
			\end{itemize}
		}
	\end{ex}
	\begin{ex}%[3P2K5-2][Loai1]
		Trên mặt chất lỏng có hai nguồn sóng kết hợp A và B cùng biên độ cùng pha cách nhau $10\ cm$. Hai điểm nguồn A và B gần như đứng yên (coi như cực tiểu dao động) và giữa chúng còn $10$ điểm đứng yên không dao động. Biết tần số rung là $26\ Hz$. Tốc độ truyền sóng là
		\choice
		{\True $47,3\ cm/s$}
		{$57,8\ m/s$}
		{$43,3\ cm/s$}
		{$27\ cm/s$}
		\loigiai{
			\begin{itemize}
				\item Coi A và B là các điểm dao động với biên độ cực tiểu $\Rightarrow$ 12 điểm dao động cực tiểu cách nhau:\\ 
				$11\dfrac{\lambda }{2} = 10\ cm$ $\Rightarrow \lambda = \dfrac{20}{11}$ cm $\Rightarrow$ $v = 47,3$ cm/s.
				\item 
				\item 
			\end{itemize}
		}
	\end{ex}
	\loai{M là cực đại, cho số cực đại giữa M và đường trung trực}
	\begin{ex}%[3P2K5-2][Loai1]
		Trong thí nghiệm giao thoa sóng trên mặt nước, hai nguồn kết hợp dao động cùng pha $O_1$ và $O_2$ cách nhau $20,5\ cm$ dao động với cùng tần số $f=15\ Hz$. Tại điểm M cách hai nguồn những khoảng $d_1 =23\ cm$ và $d_2=26,2\ cm$ sóng có biên độ cực đại. Biết rằng giữa M và đường trực của $O_1O_2$ còn một đường cực đại giao thoa. Vận tốc truyền sóng trên mặt nước là
		\choice
		{$2,4\ m/s$}
		{$16\ cm/s$}
		{$48\ cm/s$}
		{\True $24\ cm/s$}
		\loigiai{
			\begin{itemize}
				\item Tại M sóng có biên độ cực đại nên ta có: $d_1-d_2 = k \lambda$.
				\item Hai nguồn dao động cùng pha nên trung trực của hai nguồn là một cực đại giao thoa, giữa M và trung trực có một vân cực đại nên M thuộc vân cực đại $|k|=2$.
				\item Như vậy ta có $\lambda = 1,6\ cm \Rightarrow v = 24\ cm /s$.
			\end{itemize}
		}
	\end{ex}
	\begin{ex}%[3P2-6.1K][Loai1]
		Thực hiện thí nghiệm giao thoa trên mặt nước với hai nguồn kết hợp A, B khoảng cách $AB=8\ cm$, phương trình sóng tại A, B là $x_A=x_B=a\cos40\pi t\ cm$, vận tốc truyền sóng trên mặt nước là $v=30\ cm/s$. Gọi C, D là hai điểm trên mặt nước sao cho ABCD là hình vuông. Số điểm dao động với biên độ cực đại trên CD là
		\choice
		{\True $5$ điểm}
		{$11$ điểm}
		{$10$ điểm}
		{$7$ điểm}
		\loigiai{
			\begin{itemize}
				\item Ta có: $\lambda  = \dfrac{v}{f} = \dfrac{30}{20} = 1,5 $ cm.
				\item Áp dụng định lí Pitago $AC = BD = 8\sqrt{2} $.
				\item Số điểm cực đại trên CD thoả mãn:
				\begin{align*}
					\dfrac{\Delta d_C}{\lambda } \le k \le \dfrac{\Delta d_D}{\lambda} \Rightarrow \dfrac{CB-CA}{\lambda } \le k \le \dfrac{DB-DA}{\lambda} \Rightarrow \dfrac{8-8\sqrt{2}}{1,5} \le k \le \dfrac{8\sqrt{2}-8}{1,5} \Rightarrow -2,2 \le k \le 2,2.
				\end{align*}
				\item Vậy có 5 điểm dao động với biên độ cực đại trên CD.
			\end{itemize}
		}
	\end{ex}
	\begin{ex}%[3P5-3.2B][Loai5]
		Trong thí nghiệm giao thoa Young, khoảng cách hai khe là $1,2\ mm$, khoảng cách giữa mặt phẳng chứa hai khe và màn ảnh là $2\ m$. Người ta chiếu vào khe Young bằng ánh sáng đơn sắc có bước sóng $ 0,6\ \mu m $. Xét tại hai điểm M và N trên màn có tọa độ lần lượt là $6\ mm$ và $15,5\ mm$ là vị trí vân sáng hay vân tối?
		\choice
		{M sáng bậc $2$; N tối thứ $16$}
		{\True M sáng bậc $6$; N tối thứ $16$}
		{M sáng bậc $2$; N tối thứ $9$}
		{M tối bậc $2$; N tối thứ $9$}
		\loigiai{
			\begin{itemize}
				\item Ta có: $i=\dfrac{\lambda D}{a}=\dfrac{{0,6.10}^{-6}.2}{{1,2.10}^{-3}}=1\ mm\Rightarrow \left\{\begin{aligned}& \dfrac{x_m}{i}=6\Rightarrow M\ \text{là vân sáng bậc 6}\\
					& \dfrac{x}{i}=15,5\Rightarrow N\ \text{là vân tối thứ 16}
				\end{aligned}\right.$
				\item 
				\item 
			\end{itemize}
		}
	\end{ex}
	
	

\subsubsection*{PHẦN 2. Thí sinh trả lời từ câu 1 đến câu 4. Trong mỗi ý a), b), c), d) ở mỗi câu, thí sinh chọn đúng hoặc sai.}
\setcounter{ex}{0}

\begin{ex}
	Khi một sợi dây căng theo phương thẳng đứng với hai đầu cố định được kích thích dao động, hiện tượng giao thoa giữa sóng tới và sóng phản xạ tạo thành sóng dừng.
	\choiceTF[t]
	{Không tồn tại thời điểm mà sợi dây duỗi thẳng}
	{Trên dây có thể tồn tại hai điểm mà dao động tại hai điểm đó lệch pha nhau một góc là $\dfrac{\pi}{3}$}
	{Hai điểm trên dây đối xứng nhau qua một nút sóng thì dao động ngược pha nhau}
	{\True Khi giữ nguyên các điều kiện khác nhưng thả tự do đầu dưới thì không có sóng dừng ổn định} 
	\loigiai{
		\begin{itemchoice}
			\itemch Sai. Vì có tồn tại thời điểm mà sợi dây duỗi thẳng, do khoảng thời gian liên tiếp hai lần dây duỗi thẳng là nửa chu kì.
			\itemch Sai. Vì với một sợi dây với hai đầu cố định có sóng dừng. Trong các phần tử đây mà tại đó sóng tới và sóng phản xạ lệch pha nhau $\pm \dfrac{\pi}{3}+2 k \pi$ (với $k$ là các số nguyên).
			\itemch Sai. Vì hai điểm nằm trên hai bó sóng liền kề nhau đối xứng nhau qua một nút sóng thì mới dao động ngược pha nhau
			\itemch Đúng. Trên sợi dây 2 đầu cố định đang có sóng dừng nếu thả tự do đầu dưới thì sóng dừng không ổn định.
		\end{itemchoice}
	}
\end{ex}
\begin{ex}
 Trong dao động điều hòa, động năng và thế năng của vật biến thiên liên tục và chuyển hóa lẫn nhau theo quy luật. Các đại lượng này có thể được biểu diễn bằng đồ thị với hình dạng đặc trưng cho từng mối quan hệ vật lí.
	\choiceTF[t]
	{\True Đồ thị biểu diễn mối quan hệ giữa thế năng theo li độ là một đường parabol}
	{Đồ thị biểu diễn mối quan hệ giữa động năng và thế năng là một đường thẳng}
	{\True Đồ thị biểu diễn mối quan hệ giữa động năng theo li độ là một đường parabol}
	{\True Đồ thị biểu diễn mối quan hệ giữa động năng với tốc độ là một phần của đường parabol}
	\loigiai{
		\begin{itemchoice}
			\itemch Đúng.
			\itemch Sai. (là 1 đoạn thẳng không đi qua gốc tọa độ)
			\itemch Đúng.
			\itemch Đúng.
		\end{itemchoice}
	}
	%\haidong{25}
\end{ex}
\begin{ex}
	Sóng điện từ là các bức xạ lan truyền được trong chân không với tốc độ không đổi, bao gồm nhiều loại như sóng vô tuyến, tia hồng ngoại, ánh sáng nhìn thấy, tia tử ngoại, tia X và tia gamma. Mỗi loại sóng có một dải bước sóng và tần số đặc trưng. 
	\choiceTF[t]
	{Một sóng vô tuyến có tần số $10^8\ Hz$ được truyền trong không khí với tốc độ $3.10^8\ m/s$. Bước sóng của sóng đó là $1,5 \ m$}
	{Sóng điện từ có bước sóng $7.10^{-9}\ m$ thuộc tia hồng ngoại}
	{\True Sóng điện từ có bước sóng $2.10^{-8}\ m$ thuộc tia tử ngoại}
	{\True Tia X có tần số trong khoảng từ $10^{17}\ Hz$ đến $10^{19}\ Hz$}
	\loigiai{
		\begin{itemchoice}
			\itemch Sai.  $\lambda=\dfrac{v}{f}=\dfrac{3.10^{8}}{10^{8}}=3 \ m$
			\itemch Sai. Sóng điện từ có bước sóng $7.10^{-9}\ m$ thuộc tia X. Vùng tia X có bước sóng từ $10^{-9}\ m$ đến $10^{-12}\ m$
			\itemch Đúng. Sóng điện từ có bước sóng $2 . 10^{-8}\ m$ thuộc tia tử ngoại. Vùng tia tử ngoại có bước sóng từ $0,38 . 10^{-6}\ m$ đến  $10^{-9}\ m$
			\itemch Đúng. 
		\end{itemchoice}
	}
\end{ex}
\begin{ex}
	Trên mặt nước, hai nguồn đồng bộ A, B cách nhau một đoạn $30 \ cm$ có biên độ $a=1 \ cm$, tần số $f=10 \ Hz$. Biên độ sóng không đổi trong quá trình truyền đi, tốc độ truyền sóng là $v=40 \ cm/s$.
	\choiceTF[t]
	{\True Độ lớn của bước sóng là $4 \ cm$}
	{\True Trên AB có $15$ điểm là cực đại giao thoa}
	{Trên AB có $16$ điểm là cực tiểu giao thoa}
	{Trên AB, khoảng cách giữa một cực đại và cực tiểu liên tiếp là $2\ cm$}
	\loigiai{
		\begin{itemchoice}
			\itemch Đúng.
			\itemch Đúng.
			\itemch Sai.
			\itemch Sai.
		\end{itemchoice}
	}
\end{ex}
\subsubsection*{PHẦN 3. Thí sinh trả lời từ câu 1 đến câu 6.}
\setcounter{ex}{0}

\begin{ex}%[3P1K7-2][Loai1]
	Một con lắc lò xo thẳng đứng dao động điều hòa theo phương Ox, chiều dương hướng từ trên xuống. Độ cứng lò xo là $40\ N/m$. Khi qua li độ $x = 1,5\ cm$ vật chịu lực kéo đàn hồi của lò xo có độ lớn $1,6\ N$. Lấy $g = 10\ m/s^2$. Khối lượng vật nặng của con lắc là bao nhiêu kg?
	\shortans[oly]{0,1}
	\loigiai{
		\begin{itemize}
			\item Khi vật có $x = 1,5$ cm, Ox hướng xuống $\Rightarrow$ vật đang ở dưới vị trí cân bằng O đoạn 1,5cm
			\item $\Rightarrow$ Lò xo đang giãn: $\Delta \ell + 0,015$ m $\Rightarrow$ $\left|F\right| = k(\Delta \ell + 0,015) = 1,6$ N $\Rightarrow$ $\Delta \ell = 0,025$ m $\Rightarrow$ $\omega  = 20$ rad/s $\Rightarrow$ $m = 0,1$ kg.	
		\end{itemize}
	}
	%\haidong{30}
\end{ex}
\begin{ex} %[3P1B5-3][Loai1] 	
	Một vật nhỏ khối lượng $100\ g$ đang dao động điều hòa với chu kì là $0,2\ s$. Lấy $\pi^2 = 10$, mốc thế năng tại vị trí cân bằng. Tại li độ $3\sqrt{2}\ cm$, tỉ số động năng và thế năng bằng $1$. Cơ năng toàn phần của dao động là bao nhiêu mJ (làm tròn kết quả đến chữ số hàng đơn vị)?
	\shortans[oly]{180}
	\loigiai{
		\begin{itemize}
			\item Ta có: $\omega = \dfrac{2\pi}{T} = 10\pi \ rad/s\Rightarrow W_t = \dfrac{m\omega^2x^2}{2} = \dfrac{0,1.(10\pi)^2.(0,03\sqrt{2})^2}{2} = 0,09\ J$.
			\item $\dfrac{W_{\text{đ}}}{W_t} = \dfrac{W_{\text{đ}}}{0,09} = 1 \Rightarrow W_d = 0,09\ J\Rightarrow W = W_d + W_t = 0,09 + 0,09 = 0,18\ J = 180\ mJ$.
		\end{itemize}
	}
	%\haidong{30}
\end{ex}
\begin{ex}
	\immini{
		Một con lắc lò xo treo thẳng đứng gồm lò xo nhẹ và vật nhỏ A có khối lượng m. Lần lượt treo thêm các quả cân vào A thì chu kì dao động điều hòa của con lắc tương ứng là T. Hình bên biểu diễn sự phụ thuộc của $T^2$ theo tổng khối lượng $\Delta m$ của các quả cân treo vào A. Giá trị của m là bao nhiêm gam?\\
		\shortans[oly]{80}
		%	\choice
		%	{120 g}
		%	{\True 80 g}
		%	{100 g}
		%	{60 g}
	}{\begin{tikzpicture}[draw=Mapcolor,scale=0.6,transform shape]
			% Vẽ các trục tọa độ
			\draw[->] (0, 0) -- (9, 0) node[above] {\Large $\Delta m$ ($g$)};
			\draw[->] (0,0) -- (0, 6.7) node[right] {\Large $T^2$ ($s^2$)};
			%Lưới
			\draw  [help  lines, xstep=0.5cm,ystep=0.5cm,Mapcolor!50]  (0,0)  grid  (8,6);
			%Đường thẳng
			\draw[Mapcolor,line width=2pt] plot[domain=0:8] (\x,{\x/3+8/3});
			
			% Vẽ các điểm
			\filldraw (0,0) circle (4pt) node[below left=5pt] {\Large O};
			\filldraw (2,0) circle (4pt) node[below =5pt] {\Large 20};
			\filldraw (4,0) circle (4pt) node[below =5pt] {\Large 40};
			\filldraw (6,0) circle (4pt) node[below =5pt] {\Large 60};
			\filldraw (8,0) circle (4pt);
			\filldraw (1,3) circle (4pt);
			\filldraw (4,4) circle (4pt);
			\filldraw (7,5) circle (4pt);
			\filldraw (0,2) circle (4pt) node[left =5pt] {\Large 0,2};
			\filldraw (0,4) circle (4pt) node[left =5pt] {\Large 0,4};
			\filldraw (0,6) circle (4pt) node[left =5pt] {\Large 0,6};
	\end{tikzpicture}}
	\loigiai{
		\begin{itemize}
			\item Ta có: $T=2\pi\sqrt{\dfrac{m}{k}}\Rightarrow T^2=\dfrac{4\pi^2}{k}.\left(m+\Delta m\right)\Rightarrow \dfrac{0,3}{0,4}=\dfrac{m+10}{m+40}\Rightarrow m=80\ g$.
		\end{itemize}
	}	
	\haidong{15}
	
	
\end{ex}
\begin{ex}
	Trong hiện tượng giao thoa sóng trên mặt chất lỏng với hai nguồn A, B có cùng phương trình dao động $u=2 \cos 10 \pi t\  cm$, đặt cách nhau $15 \ cm$. Tốc độ truyền sóng trên mặt chất lỏng là $60 \ cm/s$. Trên AB có bao nhiêu điểm dao động với biên độ cực đại?
	\shortans[oly]{3}
	\loigiai{
		3
	}
\end{ex}
\begin{ex}
	Trong thí nghiệm Young về giao thoa ánh sáng, khoảng cách giữa hai khe là $1 \ mm$, khoảng cách từ mặt phẳng chứa hai khe đến màn quan sát là $2 \ m$. Nguồn sáng đơn sắc có bước sóng $0,45\ \mu m$. Khoảng vân giao thoa trên màn là bao nhiêu milimét?
	\shortans[oly]{0,9}
	\loigiai{
		0,9
	}
\end{ex}

\begin{ex}
	Trên mặt nước có hai nguồn sóng kết hợp A và B dao động điều hòa theo phương thẳng đứng, cùng pha, cùng tần số $40 \ Hz$. Tại điểm M trên mặt nước, cách A và B lần lượt $16 \ cm$ và $22 \ cm$, phần tử nước tại đó dao động với biên độ cực đại. Trong khoảng giữa M và đường trung trực của AB còn có $3$ đường cực đại nữa. Tốc độ truyền sóng trên mặt nước bằng bao nhiêu $cm/s$?
	\shortans[oly]{60}
	\loigiai{
		60
	}
\end{ex}